%---------------------------------------------------------------------
\section{Methods}

\subsection*{Search Strategy and Scope}

This narrative review searched PubMed and Google Scholar across five domains: (1)~\textbf{hangboard and finger-strength intervention studies} (primary literature base); (2)~\textbf{connective tissue supplementation} (collagen peptides, gelatin, omega-3, creatine, $\beta$-alanine, caffeine, magnesium); (3)~\textbf{finger and upper-limb injury management} (pulley injuries, flexor tendon strains, skin trauma); (4)~\textbf{athletic recovery modalities} (sleep, HRV, contrast water therapy, soft-tissue work, nutrition timing); (5)~\textbf{weekly training integration and periodization} (concurrent training, session sequencing, deload protocols). Searches were conducted to April 2026. Systematic reviews, meta-analyses, and practice guidelines were additionally sourced from Cochrane, GRADE working group, and Physiotherapy Evidence Database (PEDro). Practitioner claims from widely-used coaching resources (\textit{Training for Climbing}, \textit{Rock Climber's Training Manual}, Lattice Training, Dave MacLeod, Crimpd, \texttt{r/climbharder}) are reported and classified by evidence tier.

\subsection*{Inclusion and Exclusion Criteria}

\textbf{Peer-reviewed studies} were included if they: (i) used a hangboard, no-hang device, or climbing-specific resistance protocol as the primary intervention with $\geq$1 finger-strength, endurance, or tendon-structural outcome, or (ii) examined a supplementation, recovery, or rehabilitation modality with direct relevance to a climbing-specific tissue (finger flexors, annular pulleys, skin) or performance outcome. Case series and observational studies were included when RCT-level evidence was absent (injury epidemiology, diagnostic accuracy). Foundational biomechanical and physiological studies were included to contextualise training and rehabilitation rationale.

\textbf{Adjacent-sport transfer evidence} (studies in general resistance-training, running, cycling, or team-sport populations) was included only when: (a) no climbing-specific controlled study existed for the question; (b) the tissue, loading modality, or physiological mechanism was explicitly comparable; and (c) the transfer assumption was stated and rated using ESS indirectness penalties (typically ESS 3--4 vs.\ ESS 7--8 for direct climbing evidence). All claims derived from adjacent-sport evidence carry a transfer-uncertainty qualifier.

\textbf{Folklore and practitioner claims} were included when they originated from $\geq$3 independent climbing-community sources (see Claim Status Labels below). Sources with $<$3 independent attestations are classified \toverify and excluded from settled guidance.

\subsection*{Claim Status Labels}

Each claim carries one of four labels:
\begin{itemize}[noitemsep]
  \item \textbf{\evidencebacked}: peer-reviewed intervention evidence with transparent design and explicit bias context.
  \item \textbf{\bioplausible}: mechanistically plausible, indirect or incomplete intervention support; mechanism stated; direct RCT gap noted.
  \item \textbf{\folkloreverified}: community-held belief with prevalence verified from $\geq$3 independent climbing-community sources (e.g., named coaches, published books, subreddit wiki).
  \item \textbf{\toverify}: provenance unclear or fewer than 3 independent source confirmations; must not be treated as settled guidance.
\end{itemize}

\subsection*{Evidence Strength Score (ESS)}

Each actionable recommendation carries an ESS from 0--10 reflecting study design, sample relevance to climbers, risk of bias, outcome directness, and transfer indirectness:

\begin{center}
\small
\begin{tabular}{>{\raggedright}p{1.2cm} >{\raggedright\arraybackslash}p{12cm}}
\toprule
\textbf{ESS} & \textbf{Criteria} \\
\midrule
10 & Well-powered climbing RCT, low risk of bias, functional outcome (redpoint grade, route completion, injury incidence) \\
9  & Strong climbing RCT with minor limitations \\
8  & Climbing RCT with moderate limitations \\
7  & Controlled non-randomised climbing intervention or high-quality climbing systematic review / meta-analysis \\
6  & Prospective climbing cohort or repeated-measures climbing intervention with acceptable controls \\
5  & Climbing observational evidence with meaningful confounding \\
4  & High-quality adjacent-sport RCT with explicit transfer uncertainty (e.g., Achilles tendinopathy $\to$ A2 pulley) \\
3  & Observational adjacent-sport evidence with substantial indirectness \\
2  & Mechanistic rationale, expert interpretation, or in-vitro evidence; no direct intervention support \\
1  & Anecdotal, single-source, or highly limited evidence \\
0  & Unreliable, unresolvable, or critically biased source basis \\
\bottomrule
\end{tabular}
\end{center}

\noindent ESS is reduced by one level when: (i) primary outcome is a surrogate biomarker (e.g., P1NP, serum markers) rather than a functional or structural endpoint; or (ii) a study is single-source and has not been independently replicated. These reductions are documented at each recommendation site.

\subsection*{Endpoint Classification}

All recommendation sites explicitly distinguish \textbf{surrogate} outcomes (biomarkers, dynamometric measures, VAS scores, in-vitro proxies) from \textbf{functional} outcomes (on-wall climbing performance, injury incidence, return-to-sport, grade improvement). Surrogate outcomes are labelled \textit{(surrogate)} at first use; functional outcomes are labelled \textit{(functional)}. Recommendations derived solely from surrogate-endpoint studies are assigned reduced recommendation confidence regardless of ESS.

\subsection*{Synthesis Process}

Initial multi-domain literature queries used eight free-tier frontier LLMs to maximise recall across specialist topics (supplements, injury, recovery, weekly programming). Synthesis, claim classification, and ESS assignment were performed by the author with AI-assisted orchestration, with all claims reconciled against primary sources. Single-source claims carry an explicit footnote and are excluded from any recommendation presented as settled guidance. Named model brands do not appear in manuscript body text; methods may state that free-tier frontier LLMs were used for feedback and that synthesis was programmatically orchestrated. See Appendix~\ref{app:toverify} for the consolidated \toverify tracker.
